% applications/railway_vs_aircraft_dynamics-DE.tex

\chapter{Eisenbahn- und Flugzeugdynamik im Vergleich}
\label{chap:railway-vs-aircraft}

Dieses Kapitel eröffnet eine vergleichende Anwendungsperspektive, welche
die eisenbahnspezifischen Folgen geführter Bewegung verdeutlicht.

\section{Vergleichskontext}

Flugzeug- und Eisenbahndomäne teilen dynamische Prinzipien wie die
Trägheitsantwort und die Bedeutung des Komforts. Der für Anwendungen
entscheidende Unterschied ist die kinematische Freiheit: Flugbahnen
werden im freien Raum aktiv gesteuert, während Eisenbahnfahrzeuge durch
die Infrastrukturgeometrie geführt werden.

\section{Konsequenz für die Eisenbahn}

Weil die Eisenbahnbewegung durch die Infrastruktur geführt wird, werden
Alignment, Überhöhung und Laufzeitinterpretation zu den wesentlichen
Stellgrößen des Betriebsverhaltens.

Auf Anwendungsebene lautet die Ingenieurfrage häufig, wie bestehende
Geometrie durch laufzeit- und realisierungsbewusste Interpretation
zulässige Betriebsbereiche unterstützen kann.

\section{Interpretation in AIM}

Innerhalb von AIM begründet dieser Vergleich die Verwendung von
kanonischem pose3, Fahrzeugrauminterpretation und betrieblicher
Laufzeitauswertung für die Entscheidungsunterstützung im Eisenbahnwesen.

Der Vergleich dient ausschließlich der Erklärung; er führt kein
paralleles Architekturmodell ein.

\section{Verbindung zu AIM}

\begin{summary}
Der Vergleich von Eisenbahn und Flugzeug verdeutlicht, warum
Eisenbahnanwendungen eine starke Kopplung von Infrastrukturseman-tik,
Laufzeitdiensten und betrieblicher Interpretation erfordern.

Er bekräftigt die Anwendungsfolgen kanonischer AIM-Begriffe, ohne diese
neu zu definieren.
\end{summary}
